
\documentclass[12pt]{article}
\usepackage[utf8]{inputenc}
\usepackage{amsmath,amssymb}
\usepackage{bm}

\title{Ecuaciones Básicas}
\date{}

\begin{document}

\maketitle

\section*{1 Ecuaciones Básicas}

\subsection*{1.1 Ecuaciones gobernantes}

Sea $(0,T) \subset \mathbb{R}$ un intervalo de tiempo, durante el cual seguimos el movimiento del fluido, y sea $\Omega \subset \mathbb{R}^N$, $N=1,2,3$, el dominio ocupado por el fluido. (Por simplicidad, suponemos que es independiente del tiempo $t$).

\subsubsection*{1.1.1 Descripción del flujo}

Existen dos posibilidades para describir el movimiento del fluido: Lagrangiana y Euleriana.

Aquí utilizaremos la descripción Euleriana basada en la determinación de la velocidad $\bm{v}(x,t) = (v_1(x,t),\dots,v_N(x,t))$ de la partícula de fluido que pasa por el punto $x$ en el instante $t$.

Usaremos la siguiente notación básica: $\rho$ — densidad del fluido, $p$ — presión, $\theta$ — temperatura absoluta. Estas cantidades se denominan variables de estado.

A continuación, introduciremos la formulación matemática de las leyes físicas fundamentales: la ley de conservación de la masa, la ley de conservación del momento lineal y la ley de conservación de la energía, llamadas en conjunto leyes de conservación.

\subsubsection*{1.1.2 La ecuación de continuidad}

\begin{equation}
\frac{\partial \rho}{\partial t}(x,t) + \operatorname{div}(\rho(x,t)\bm{v}(x,t)) = 0, \quad t \in (0,T), \ x \in \Omega,
\tag{1.1.1}
\end{equation}

es la forma diferencial de la ley de conservación de la masa.

\subsubsection*{1.1.3 Las ecuaciones del movimiento}

Las ecuaciones dinámicas básicas equivalentes a la ley de conservación del momento lineal tienen la forma

\begin{equation}
\frac{\partial}{\partial t}(\rho \bm{v}_i) + \operatorname{div}(\rho \bm{v}_i v) = \rho \bm{f}_i + \sum_{j=1}^{N} \frac{\partial \tau_{ji}}{\partial x_j}, \quad i = 1,\dots,N.
\tag{1.1.2}
\end{equation}

Esto se puede escribir como

\begin{equation}
\frac{\partial}{\partial t}(\rho v) + \operatorname{div}(\rho \bm{v} \otimes \bm{v}) = \rho \bm{f} + \operatorname{div}\mathbb{T}.
\tag{1.1.3}
\end{equation}

Aquí $\bm{v} \otimes \bm{v}$ es el tensor con componentes $\bm{v}_i v_j$, $i,j = 1,\dots,N$, $\tau_{ij}$ son los componentes del tensor de esfuerzos $\mathbb{T}$ y $\bm{f}_i$ son los componentes de la densidad de fuerza volumétrica externa $f$.

\subsubsection*{1.1.4 La ley de conservación del momento de momento angular}

\textbf{Teorema 1.1} La ley de conservación del momento de momento angular es válida si y sólo si el tensor de esfuerzos $\mathbb{T}$ es simétrico.

\subsubsection*{1.1.5 Las ecuaciones de Navier–Stokes}

Las relaciones entre el tensor de esfuerzos y otras cantidades que describen el flujo del fluido, en particular la velocidad y sus derivadas, representan las llamadas ecuaciones reológicas del fluido. La ecuación reológica más simple

\begin{equation}
\mathbb{T} = -p\mathbb{I},
\tag{1.1.4}
\end{equation}

caracteriza a un fluido incompresible ideal. Aquí $p$ es la presión e $\mathbb{I}$ es el tensor identidad:

\begin{equation}
\mathbb{I} = 
\begin{pmatrix}
1 & 0 & 0 \\
0 & 1 & 0 \\
0 & 0 & 1
\end{pmatrix}, \quad \text{para } N = 3.
\tag{1.1.5}
\end{equation}

Además de las fuerzas de presión, en los fluidos reales también actúan fuerzas de corte por fricción como consecuencia de la viscosidad. Por tanto, en el caso de un fluido viscoso, se añade una contribución $\mathbb{T}'$ que caracteriza el esfuerzo cortante al término $-p\mathbb{I}$:

\begin{equation}
\mathbb{T} = -p\mathbb{I} + \mathbb{T}'.
\tag{1.1.6}
\end{equation}

\end{document}


\subsubsection*{1.1.6 Propiedades de los coeficientes de viscosidad}

Aquí $\mu$ y $\lambda$ se llaman primer y segundo coeficiente de viscosidad, respectivamente. El coeficiente $\mu$ también se denomina viscosidad dinámica. A partir de la teoría cinética de gases se derivan las condiciones

\begin{equation}
\mu \geq 0, \qquad 3\lambda + 2\mu \geq 0.
\tag{1.1.7}
\end{equation}

Para gases monoatómicos se cumple

\begin{equation}
3\lambda + 2\mu = 0.
\tag{1.1.8}
\end{equation}

Esta condición se usa generalmente incluso en el caso de gases más complejos.

\subsubsection*{1.1.7 La ecuación de la energía}

La ley de conservación de la energía se expresa en la ecuación de la energía. Para este propósito, definimos la energía total

\begin{equation}
E = \rho\left(e + \frac{|\bm{v}|^2}{2}\right),
\tag{1.1.9}
\end{equation}

donde $e$ es la energía interna específica. La ecuación de energía tiene la forma

\begin{equation}
\frac{\partial E}{\partial t} + \operatorname{div}(E \bm{v}) = \rho \bm{f} \cdot \bm{v} + \operatorname{div}(\mathbb{T} \bm{v}) + \rho q - \operatorname{div} q.
\tag{1.1.10}
\end{equation}

Para un fluido newtoniano, esta se convierte en

\begin{equation}
\frac{\partial E}{\partial t} + \operatorname{div}(E \bm{v}) = \rho \bm{f} \cdot \bm{v} - \operatorname{div}(p \bm{v}) + \operatorname{div}(\lambda \bm{v} \operatorname{div} \bm{v}) + \operatorname{div}(2\mu D(\bm{v}) \bm{v}) + \rho q - \operatorname{div} q,
\tag{1.1.11}
\end{equation}

donde $q$ es la densidad de fuentes térmicas, y $q$ también denota el flujo de calor, el cual se determina mediante la temperatura de acuerdo a la ley de Fourier:

\begin{equation}
q = -k \nabla \theta,
\tag{1.1.12}
\end{equation}

donde $k \geq 0$ es el coeficiente de conductividad térmica.

\subsection*{1.2 Relaciones termodinámicas}

Para completar el sistema de leyes de conservación, deben añadirse ecuaciones adicionales provenientes de la termodinámica.

Las cantidades $\theta$, $\rho$ y $p$ se denominan variables de estado. Todas son funciones positivas. El gas se caracteriza por la ecuación de estado

\begin{equation}
p = p(\rho, \theta)
\tag{1.2.1}
\end{equation}

y por la relación

\begin{equation}
e = e(\rho, \theta).
\tag{1.2.2}
\end{equation}

Consideramos el llamado gas perfecto (o ideal), cuyas variables de estado satisfacen la ecuación

\begin{equation}
p = R \theta \rho,
\tag{1.2.3}
\end{equation}

donde $R > 0$ es la constante del gas, que puede expresarse como

\begin{equation}
R = c_p - c_v,
\tag{1.2.4}
\end{equation}

donde $c_p$ y $c_v$ son los calores específicos a presión y volumen constantes, respectivamente. Los experimentos muestran que $c_p > c_v$, y por tanto $R > 0$, y que $c_p$ y $c_v$ pueden considerarse constantes en un amplio rango de temperaturas. Se define

\begin{equation}
\gamma = \frac{c_p}{c_v} > 1,
\tag{1.2.5}
\end{equation}

que se denomina constante adiabática de Poisson. Por ejemplo, para el aire seco $\gamma = 1.4$. La energía interna específica del gas perfecto se da por

\begin{equation}
e = c_v \theta.
\tag{1.2.6}
\end{equation}

\subsubsection*{1.2.1 Entropía}

Una cantidad termodinámica importante es la entropía $S$, definida por la relación

\begin{equation}
\theta dS = de + p dV,
\tag{1.2.7}
\end{equation}

donde $V = 1/\rho$ es el llamado volumen específico. Esta identidad se deduce de la termodinámica, bajo la suposición de que la energía interna es una función de $S$ y $V$: $e = e(S, V)$, lo cual explica el significado de los diferenciales.

\textbf{Teorema 1.3.} Para un gas perfecto,

\begin{equation}
S = c_v \ln \left( \frac{p/p_0}{(\rho/\rho_0)^\gamma} \right) + \text{constante},
\tag{1.2.8}
\end{equation}

donde $p_0$ y $\rho_0$ son valores de referencia de presión y densidad.

\subsubsection*{1.2.2 Flujo barotrópico}

Decimos que el flujo es barotrópico si la presión puede expresarse como función de la densidad:

\begin{equation}
p = p(\rho),
\tag{1.2.9}
\end{equation}

lo cual implica que $p(x,t) = p(\rho(x,t))$, es decir, $p = p \circ \rho$. Suponemos que

\[
p : (0, +\infty) \to (0, +\infty)
\]

y que es una función continua y diferenciable.

